\subsubsection{Spike-count reconstruction on held-out macaque trials}
\label{subsubsection:rebuttal_reconstruction}

To compare how much of the recorded activity each method's representation retains, we reconstruct spike counts from the representations compared on these data (the 192-dimensional NLDisco level-192 latents of seed 1, CEBRA-Time, and LangevinFlow, and the 191-dimensional ST-NDT representation, whose width the model ties to the number of channels), using the saved representations and checkpoints without retraining. The target of each window is the vector of spike counts of all 191 channels in its final bin, the bin its representation is aligned to.

\textit{Shared readout.} For every method we fit the same Poisson regression with a log link and an L2 penalty, from the full representation (each dimension standardized with training statistics) to the final-bin counts of all channels. It is fit on the training trials of the whole-trial 80/10/10 split, its penalty strength is chosen from $\{10^{-6},10^{-5},\dots,10^{0}\}$ by bits per spike on the validation trials, and it is scored once on the test trials (13,020 windows from 191 trials).

\textit{Native reconstruction.} NLDisco, LangevinFlow, and ST-NDT also reconstruct the final bin with their own decoders: NLDisco's linear temporal decoder, with the channel standardization undone; LangevinFlow's Transformer/Poisson decoder, with dropout disabled and posterior means; and ST-NDT's Poisson output, run on the unmasked window in evaluation mode. For all three, as for the shared readout, the input window contains the reconstructed bin. NLDisco's decoder is trained with an MSE loss and an identity output, so 0.5\% of its predictions are nonpositive; its predictions are floored at 0.1 counts per bin, a value chosen by validation bits per spike. CEBRA-Time has no decoder.

\textit{Scores.} Bits per spike is the reconstruction score of the Neural Latents Benchmark~\citep{pei_2022_nlb}: the Poisson log-likelihood of the true counts under the predicted rates minus that under each channel's mean test rate, divided by the number of test spikes and expressed in bits. Variance explained is $1-\sum_t (y_{tc}-\hat{y}_{tc})^2/\sum_t (y_{tc}-\bar{y}_c)^2$ for each channel $c$, averaged over channels. For both, zero corresponds to predicting each channel's mean rate and higher is better. Intervals are 95\% bootstrap intervals over whole test trials (1,000 resamples).

With the shared readout, ST-NDT's representation reconstructs held-out spike counts best (0.253 bits per spike and 0.515 variance explained) (\autoref{table:rebuttal_reconstruction}). NLDisco and LangevinFlow reconstruct them equally well within the confidence intervals (bits per spike 0.156 and 0.159; variance explained 0.335 for both), and both are above CEBRA-Time (0.115 and 0.235). With each method's own decoder, NLDisco reconstructs best (0.155 bits per spike and 0.340 variance explained, compared with 0.148 and 0.314 for ST-NDT and 0.141 and 0.296 for LangevinFlow). NLDisco's 192 sparse latents thus retain as much information about the recorded activity as LangevinFlow's dense representation of the same width.

\begin{table}[!htbp]
    \centering
    {\small\setlength{\tabcolsep}{4pt}
    \begin{tabular}{lcccc}
\toprule
& \multicolumn{2}{c}{Shared Poisson readout} & \multicolumn{2}{c}{Native decoder} \\
\cmidrule(lr){2-3}\cmidrule(lr){4-5}
Method & Bits/spike & Var.\ explained & Bits/spike & Var.\ explained \\
\midrule
NLDisco & 0.156 {\scriptsize [0.153, 0.157]} & \textbf{0.335} {\scriptsize [0.331, 0.338]} & \textbf{0.155} {\scriptsize [0.153, 0.157]} & \textbf{0.340} {\scriptsize [0.336, 0.343]} \\
CEBRA-Time & 0.115 {\scriptsize [0.113, 0.117]} & 0.235 {\scriptsize [0.231, 0.238]} & -- & -- \\
LangevinFlow & \textbf{0.159} {\scriptsize [0.157, 0.161]} & \textbf{0.335} {\scriptsize [0.323, 0.341]} & 0.141 {\scriptsize [0.139, 0.142]} & 0.296 {\scriptsize [0.293, 0.299]} \\
\bottomrule
\end{tabular}
}
    \caption{\textbf{Final-bin spike-count reconstruction on held-out test trials.} Each representation (192-dimensional; 191 for ST-NDT) predicts the counts of all 191 channels in the bin it is aligned to. The shared readout is the same validation-tuned L2 Poisson regression for every method; the native columns use each model's own decoder. Brackets: 95\% whole-trial bootstrap intervals. Column maxima at the displayed precision are bolded.}
    \label{table:rebuttal_reconstruction}
\end{table}
